% ch3_answers.tex -- answers to the Quick check questions in Chapter 3 (Module 6)
\begin{chapteranswers}{3}{Electromagnetic Induction: Flux, Faraday and Lenz}
  % ---------- 3.1 Magnetic flux
  \item $r = \qty{3.0}{cm} = \qty{0.030}{m}$, so
        $A = \pi r^2 = \pi \times (\qty{0.030}{m})^2 = \qty{2.827e-3}{m^2}$.
        The field is along the normal, so $\theta = \ang{0}$.\\[4pt]
        $\Phi = BA\cos\theta = \qty{0.20}{T} \times \qty{2.827e-3}{m^2} \times 1
        = \qty{5.65e-4}{Wb}$. \ Answer: \final{$\Phi = \qty{5.7e-4}{Wb}$}.
  \item The angle to the plane is \ang{25}, so $\theta = \ang{90} - \ang{25} = \ang{65}$.\\[4pt]
        $\Phi = \qty{0.20}{T} \times \qty{2.827e-3}{m^2} \times \cos\ang{65}$\\[4pt]
        $\phantom{\Phi} = \qty{5.65e-4}{Wb} \times 0.423 = \qty{2.39e-4}{Wb}$.
        \ Answer: \final{$\Phi = \qty{2.4e-4}{Wb}$}.
  \item Magnetic flux $\Phi$ measures how much field goes through a whole loop or
        surface. Its unit is the weber (\si{Wb}). Magnetic flux density $B$ is
        the field strength: the flux through each square metre that faces the
        field. Its unit is the tesla (\si{T} $= \si{Wb.m^{-2}}$). Flux depends on
        the area and the angle, whereas flux density does not.
  \item (a) Square-on means $\theta = \ang{0}$.
        $\Phi_\mathrm{X} = \qty{0.20}{T} \times \qty{0.050}{m^2} \times \cos\ang{0}
        = \final{\qty{0.010}{Wb}}$; \
        $\Phi_\mathrm{Y} = \qty{0.50}{T} \times \qty{0.010}{m^2} \times 1
        = \final{\qty{0.0050}{Wb}}$.\\[4pt]
        (b) The claim is \textbf{incorrect}. X has twice the flux of Y, but X's
        field (\qty{0.20}{T}) is weaker than Y's (\qty{0.50}{T}). Flux depends on
        $B$, $A$ and $\theta$. X has more flux because its loop is five times
        bigger. More flux means a stronger field only when the areas and angles
        are the same.
  % ---------- 3.2 Faraday's law
  \item $\Delta\Phi = \qty{1.5e-3}{Wb} - \qty{6.0e-3}{Wb} = \qty{-4.5e-3}{Wb}$.\\[4pt]
        $\lvert\varepsilon\rvert = N\dfrac{\lvert\Delta\Phi\rvert}{\Delta t}
        = \dfrac{120 \times \qty{4.5e-3}{Wb}}{\qty{0.30}{s}} = \qty{1.8}{V}$.
        \ Answer: \final{$\varepsilon = \qty{1.8}{V}$}.
  \item $A = \qty{25e-4}{m^2} = \qty{2.5e-3}{m^2}$. The flux falls from
        $BA = \qty{0.15}{T} \times \qty{2.5e-3}{m^2} = \qty{3.75e-4}{Wb}$ to zero.\\[4pt]
        $\lvert\varepsilon\rvert = \dfrac{40 \times \qty{3.75e-4}{Wb}}{\qty{0.050}{s}}
        = \qty{0.30}{V}$; \quad
        $I = \dfrac{\varepsilon}{R} = \dfrac{\qty{0.30}{V}}{\qty{2.5}{\ohm}} = \qty{0.12}{A}$.\\[4pt]
        Answer: \final{$\varepsilon = \qty{0.30}{V}$, \ $I = \qty{0.12}{A}$}.
  \item Rearrange Faraday's law for the time:
        $\Delta t = \dfrac{N\lvert\Delta\Phi\rvert}{\lvert\varepsilon\rvert}
        = \dfrac{500 \times \qty{8.0e-4}{Wb}}{\qty{2.0}{V}} = \qty{0.20}{s}$.
        A longer time gives a smaller emf, so this is the longest time.
        \ Answer: \final{$\Delta t = \qty{0.20}{s}$}.
  \item \textbf{Yes}, an emf is induced: the flux through the coil changes, and
        Faraday's law needs only a change in flux. \textbf{No}, no current flows:
        the circuit is open, so charges cannot flow around it. The emf appears as a
        potential difference between the two loose ends.
  \item The emf depends on the \emph{rate of change} of flux,
        $\varepsilon = -N\dfrac{\Delta\Phi}{\Delta t}$, not on the size of the flux.
        A steady flux has $\Delta\Phi = 0$, so $\varepsilon = 0$, however big the
        flux is.
  % ---------- 3.3 Lenz's law
  \item The induced current flows in the direction that makes a magnetic field
        opposing the change in flux that produced it. (For an open circuit: the
        induced emf acts in the direction that would drive a current whose
        magnetic field opposes the change in flux that produced it.)
  \item Grip rule for the wire: thumb to the right (along $I$). Above the wire your
        fingers come out of the page, so the field at the loop points out of the
        page (dots). By $B = \dfrac{\mu_0 I}{2\pi r}$ the field gets weaker as $r$
        grows. So as the loop moves away, the out-of-page flux decreases (fewer
        dots). By Lenz's law the induced field opposes this decrease, so it points
        out of the page inside the loop. The grip rule for a loop then gives an
        \final{anticlockwise} current, as seen on the page.
  \item The magnet is moving away, so the coil attracts it: the near face becomes
        an \final{N pole} (N attracts S). An N face has an \final{anticlockwise}
        current when you look at it, so the current is anticlockwise viewed from
        the magnet. (Flux check: field lines go into an S pole, so the field
        through the coil points towards the magnet. This flux is decreasing, so
        the induced field also points towards the magnet: out of the near face,
        which is an N face.)
  \item The area decreases, so the into-page flux decreases (fewer crosses). The
        induced field opposes the decrease, so it points into the page. The
        right-hand grip rule gives a \final{clockwise} current.
  \item (1) Lenz's law: the induced current opposes the change in flux that causes
        it, so a coil must \emph{repel} a magnet pushed towards it. (2) Because of
        this opposing force, work must be done to push the magnet in. This work
        becomes electrical energy and then heat, so energy is only changed in
        form. (3) If the coil attracted the magnet, a small push would make the
        magnet speed itself up, giving more current and a stronger pull.
        Kinetic and electrical energy would come from nothing, which breaks the
        law of conservation of energy. So the student's coil is impossible.
  % ---------- 3.4 Motional emf
  \item The wing is a rod moving across the vertical part of the field:
        $\varepsilon = lvB = \qty{32}{m} \times \qty{240}{m.s^{-1}} \times \qty{5.0e-5}{T}
        = \qty{0.384}{V}$. \ Answer: \final{$\varepsilon = \qty{0.38}{V}$}.
  \item \final{Zero.} The velocity is parallel to $B$, so the force on each charge is
        $F = qvB\sin\ang{0} = 0$. No charges are pushed along the rod.
        (Faraday's view: the rod moves along the field lines, not across them.
        The area it sweeps out is parallel to $B$, so no extra flux passes
        through the circuit.)
  \item Palm rule for a positive charge: fingers out of the page (along $B$), thumb
        up the page (along $v$). The palm pushes to the right. So positive charge
        collects at the \final{right-hand end}.
  \item $\varepsilon = IR = \qty{0.15}{A} \times \qty{2.0}{\ohm} = \qty{0.30}{V}$.\\[4pt]
        $v = \dfrac{\varepsilon}{lB} = \dfrac{\qty{0.30}{V}}{\qty{0.50}{m} \times \qty{0.40}{T}}
        = \qty{1.5}{m.s^{-1}}$.\\[4pt]
        The magnetic force on the rod is
        $F = lIB = \qty{0.50}{m} \times \qty{0.15}{A} \times \qty{0.40}{T} = \qty{0.030}{N}$,
        against the motion. At a steady speed the push must balance it.\\[4pt]
        Answer: \final{$v = \qty{1.5}{m.s^{-1}}$, \ push $= \qty{0.030}{N}$} in the
        direction of motion.
  % ---------- 3.5 Solenoids, two coils and metal plates
  \item Switch closed: a brief flick one way, because the primary current and its
        flux grow. Kept closed: zero, because the current and the flux are steady.
        Switch opened: a brief flick the \emph{other} way, because the flux now
        decreases (Lenz's law gives the opposite direction).
  \item Pushed in: a brief flick, because the flux through the large coil
        increases (Faraday). Held still: zero, because the flux is steady. Pulled
        out: a brief flick the \emph{other} way, because the flux decreases, and
        Lenz's law gives the opposite direction. The steady current alone
        induces nothing; only the change in flux caused by the motion does.
  \item An emf is induced only while the flux through the secondary changes. AC
        changes size and direction all the time, so the primary's field and the
        flux through the secondary change all the time. This gives a continuous
        (alternating) emf. A steady DC current gives a steady flux, so there is no
        emf, except for brief pulses when it is switched on or off.
  \item The \final{solid plate} stops first. As it swings into and out of the field,
        the flux through parts of it changes, so large eddy currents are induced
        (Faraday). By Lenz's law, the forces on these currents oppose the motion,
        and the kinetic energy becomes heat. In the slotted plate the slots break
        the current loops, so the eddy currents and the braking force are much
        smaller.
  \item Field at the start:
        $B = \dfrac{\mu_0 NI}{L} = \dfrac{4\pi\times10^{-7}\,\si{N.A^{-2}} \times 800 \times \qty{3.0}{A}}{\qty{0.40}{m}}
        = \qty{7.54e-3}{T}$.\\[4pt]
        $A = \qty{5.0}{cm^2} = 5.0 \times 10^{-4}\ \si{m^2}$. The coil is wound
        tightly, so each of its turns has this area and all of the solenoid's flux
        passes through it.\\[4pt]
        Change in flux (from $BA$ to zero):
        $\lvert\Delta\Phi\rvert = \qty{7.54e-3}{T} \times \qty{5.0e-4}{m^2} = \qty{3.77e-6}{Wb}$.\\[4pt]
        $\lvert\varepsilon\rvert = \dfrac{25 \times \qty{3.77e-6}{Wb}}{\qty{0.025}{s}}
        = \qty{3.77e-3}{V}$. \ Answer: \final{$\varepsilon = \qty{3.8e-3}{V}$} (\qty{3.8}{mV}).
\end{chapteranswers}
